\subsection{Proof Construction and Revision}

A proof is a checkable argument from declared hypotheses, definitions, and earlier proved results to the exact stated conclusion.

\begin{proofmethod} \label{meth:proof-construction-audit} \hlt{Proof Construction and Audit}\\
Given a theorem-style statement, construct and review its proof as follows.
\begin{enumerate}[label=(\roman*)]
\setlength{\itemsep}{0pt}
\item \hlt{Parse}: expose the objects, hypotheses, quantifier dependencies, and conclusion by Proof Method \ref{meth:proof-grammar}.
\item \hlt{Plan}: select a direct argument, cases, two inclusions, induction, contradiction, construction, or another fitting architecture.
\item \hlt{Dependencies}: list only earlier results whose hypotheses can be verified at the point of citation.
\item \hlt{Obligations}: map each numbered claim, implication direction, existence witness, and uniqueness assertion to a proof unit.
\item \hlt{Execution}: introduce variables before use and justify each inference by a definition, calculation, hypothesis, or earlier result.
\item \hlt{Boundaries}: handle empty, zero, singular, endpoint, and other degenerate cases that affect an argument.
\item \hlt{Closure}: verify that the final sentence discharges the original conclusion rather than a nearby surrogate.
\item \hlt{Cold read}: reread without relying on memory; test notation, citations, and transitions as a skeptical reader would.
\end{enumerate}
The proof should expose its plan when the route is not apparent, yet retain every step needed for verification. Mathematical-writing guidance likewise emphasizes status, motivation, and conceptual structure \citep{mit1998mathematicalwriting,knuth1988mathematicalwriting}.
\end{proofmethod}

\begin{process} \label{proc:proof-revision-passes} \hlt{Proof Revision Passes}\\
Revise a completed proof in three passes.
\begin{enumerate}[label=(\roman*)]
\setlength{\itemsep}{0pt}
\item \hlt{Validity}: check every dependency, implication, case split, witness, index, limit, and exceptional case.
\item \hlt{Architecture}: align proof units with statement claims and make long operative stages visible.
\item \hlt{Exposition}: remove repetition, bind notation earlier, compress routine algebra, and expand decisive unfamiliar steps.
\end{enumerate}
Typography is reviewed only after the mathematical and logical structure is stable.
\end{process}

\begin{warning} \label{warn:evidence-is-not-proof} \hlt{Evidence Is Not Proof}\\
Finite testing, a diagram, numerical agreement, or a successful program run supports only the cases and conditions examined. None closes an unrestricted quantified claim without an exhaustive argument and a proof that the enumeration is complete. Computation may discover a conjecture, locate a counterexample, or verify a formally specified finite case.
\end{warning}

\begin{principle} \label{princ:proof-detail-proportionality} \hlt{Proof Detail Proportionality}\\
Detail belongs where correctness or understanding depends on it. Expand novel constructions, fragile estimates, changes of quantifier, and boundary arguments; compress familiar algebra only while the remaining chain is checkable. The goal is neither maximal formality nor persuasive rhetoric, but a reproducible mathematical argument.
\end{principle}
